% TeX root = ../Main.tex

% First argument to \section is the title that will go in the table of contents. Second argument is the title that will be printed on the page.
\section[Question -- {\it Derivatives}]{}

\subsection{Derivative in Computer Vision}
Image derivatives are the mathematical foundation of both edge and corner detection, as they capture local intensity variations that correspond to meaningful structural features.
A 2D image can be treated as 2D discrete function $I(x,y)$. The derivatives are approximated via finite
differences or convolution with derivative kernels.

The first order derivatives explains the rate of change of intensity in an image, representing the edges as local maximum $\nabla I (\frac{\partial I}{\partial x}),\frac{\partial I}{\partial y}$
The second order derivative describes the rate of change of the cradient, and the edges are represented as zero crossing point (as in LoG) $\nabla^2 I (\frac{\partial^2 I}{\partial x^2}),\frac{\partial^2 I}{\partial y^2}$
However the LoG function is high computational complex, and it is approximated by DoG for computational efficiency.

\subsection{Gradient magnitude and direction}
From the first-order derivatives $I_x$ and $I_y$, — typically computed via the Sobel operator, a separable filter combining a derivative kernel with a smoothing kernel — 
two quantities are derived:
\begin{itemize}
\item Magnitude: $|\nabla I| = \sqrt{I_x ^2 + I_y ^2}$, which encodes the edge strenght (high values represents a sharp transition)
\item Direction: $\theta = arctan(\frac{I_y}{I_x})$, which encodes the edge orientation (perpendicular to edge contours)
\end{itemize}
These quantities are crucial to apply NMS, since the edge direction is used to determine the orientation of the local neighborhood for non-maximum suppression, ensuring that only the most significant edge pixels are retained 
and edge magnitude is used to threshold the edges, discarding weak edges that are likely to be noise.

\subsection{Second order derivative for Intensity transitions}
The laplacian (LoG) responds to intensity transitions with a zero corssing at the edge location, remarking edge boundary with high spatial accuracy.
However, its main challenge is noise since it is amplified by the function, solved by a pre-smoothing by a Gaussian filter.
Two main strategies are used to overcome the noise sensitivity:
\begin{itemize}
    \item Gaussian smoothing, before differentiation
    \item Medial filtering, non linear strategy for salt and pepper noise
\end{itemize}

\subsection{Corner detection}
Corner detection relies on image derivatives and a second-moment matrix approach (structure tensor). The second moment matrix $H$ over a window $W$ is defined as:
\[
  H = \sum_{(x,y) \in W} \begin{bmatrix} I_x^2 & I_x I_y \\ I_x I_y & I_y^2 \end{bmatrix}
\]

To analyze the local intensity variation, we study the eigenvalues $\lambda_1$ and $\lambda_2$ of $H$, which correspond to the magnitudes of maximum and minimum intensity change:
\begin{itemize}
    \item $\lambda_1 \approx \lambda_2 \approx 0$ (Flat region): The error is low in all directions; intensity is nearly constant.
    \item $\lambda_1 \gg \lambda_2 \approx 0$ (Edge): The error is high in one direction but low along the edge.
    \item $\lambda_1 \approx \lambda_2 \gg 0$ (Corner): The error is high in all directions, representing a point of distinct intersection.
\end{itemize}

Since calculating eigenvalues for every pixel is computationally expensive, the Harris detector introduces a response function $R$ that approximates their behavior without explicit eigenvalue decomposition:
\[
  R = \det(H) - k \cdot (\operatorname{trace}(H))^2 = \lambda_1 \lambda_2 - k(\lambda_1 + \lambda_2)^2
\]
where $k$ is an empirical constant (typically $0.04 - 0.06$). The value of $R$ indicates the "cornerness":
\begin{itemize}
    \item $R$ large and positive $\implies$ Corner.
    \item $R$ large and negative $\implies$ Edge.
    \item $|R|$ small $\implies$ Flat region.
\end{itemize}

Properties of the Harris detector:
\begin{itemize}
    \item \textbf{Equivariant} to translation and rotation (eigenvalues and shape are preserved under rotation).
    \item \textbf{Invariant} to additive brightness shifts ($I' = I + b$) because derivatives eliminate constant offsets.
    \item \textbf{Not scale-invariant}, as a physical corner may appear as an edge or a flat region when zoomed in.
\end{itemize}